\section{Collected Dataset}
\subsection{Motivation}
The personalized DNS dataset was extensively gathered, but its training samples did not accurately represent the problem we aimed to solve. The DNS dataset prioritized quantity over quality, a common strategy in deep learning for rapid and cost-effective data collection.
Due to limited computing resources, we initially trained our models using only 100 hours of data from the 500 hours available in the personalized DNS dataset. Unfortunately, the results of these experiments were unsatisfactory, as the dataset lacked both quality and quantity.

\subsection{Data Collection}
The dataset replicates the scenario of two speakers, Speaker A and Speaker B, speaking simultaneously and provides a separate clean reference speech for Speaker A.
\subsubsection{Train-validation split}
The model was trained and validated using a train-validation split, consisting of 30 speakers with 10 clean 10-second utterances each. Among the speakers, 22 were male and 8 were female, all speaking Arabic. For each speaker, pairs of speaker-reference and silence-reference were generated, along with 10 samples from other speakers for speaker B. The total number of samples in the train-validation split is 30000. The validation set represents 1\% of the train-validation split, totaling 300 samples.

\subsubsection{Test split}
The test split was used to assess the model's performance with new speakers. It comprises 6 speakers, each providing 5 clear 10-second utterances. These individuals include 2 male Arabic speakers, 3 female Arabic speakers, and 1 female English speaker. This distribution was deliberately chosen to gauge the model's ability to generalize to different speakers. The generation process is akin to the train-validation split, with the exception of the inclusion of silence samples. The test split consists of a total of 1200 samples.